\section{Normalización constitutiva e incidencias espectrales del retorno}
\label{iii:sec:incidencias-espectrales}

La respuesta de las dos hojas y el funcional de retorno conservan dos
clases complementarias de información. La primera comprende el producto
y el cociente de los autovalores constitutivos; la segunda comprende las
alturas, los períodos y las multiplicidades orientadas de las incidencias.
Esta sección reúne las operaciones que permiten distinguir esas lecturas.
Su antecedente es el transporte angular producido por APP--TRIT--TPK,
con sus proyectores y su coordenatización levantada. Las normalizaciones se aplican
después de construir ese transporte y de conservar su orientación.

\subsection{Separación exacta del modo común y del modo orientado}

Sean $r_+,r_->0$ los autovalores etiquetados de
\eqref{iii:eq:rchannels}. Definamos
\begin{equation}
 a=\sqrt{r_+r_-},\qquad
 \eta=\frac12\log\frac{r_+}{r_-},\qquad
 \widetilde{\mathcal V}=a^{-1}\mathcal V.
 \label{iii:eq:rec27-normalizacion}
\end{equation}
En la cámara $x>y>0$ se tiene $a\in(3/4,1)$ y $\eta>0$.
La marca $\mathcal R$ conserva $\mathcal R P_\pm=\pm P_\pm$.

\begin{proposition}[Factorización positiva con orientación conservada]
\label{iii:prop:rec27-factorizacion}
La respuesta posee la factorización única
\begin{equation}
 \mathcal V=a\exp(\eta\mathcal R),\qquad
 \det\widetilde{\mathcal V}=1.
 \label{iii:eq:rec27-unimodular}
\end{equation}
Sus cuatro lecturas constitutivas satisfacen
\begin{equation}
 \widehat\varepsilon=a^2e^{2\eta},\quad
 \widehat\mu=a^2e^{-2\eta},\quad
 \widehat Z=e^{-2\eta},\quad
 \widehat c=a^{-2}.
 \label{iii:eq:rec27-cuatro}
\end{equation}
Con la orientación fija, el intercambio de hojas conserva $a$ y cambia
$\eta$ por $-\eta$.
\end{proposition}
\begin{proof}
El cálculo espectral da
$\exp(\eta\mathcal R)=e^\eta P_++e^{-\eta}P_-$.
Las definiciones implican $ae^\eta=r_+$ y $ae^{-\eta}=r_-$,
de modo que se recupera exactamente $\mathcal V$.
El producto de sus autovalores fija la única raíz positiva $a$;
su cociente fija el único valor real de $\eta$.
La sustitución en \eqref{iii:eq:four} prueba las cuatro identidades.
Finalmente, intercambiar los autovalores invierte su cociente y conserva
su producto.
\end{proof}

La normalización unimodular conserva toda la coordenada orientada y
separa la coordenada común. Si se conserva también $a$, la transformación
es reversible. Si se expresa únicamente $\widetilde{\mathcal V}$,
queda determinada $\widehat Z$, mientras $\widehat c$ necesita el
factor común archivado. Esta distinción explica el papel de las hojas
recíprocas en un lector de producto unitario: tal lector realiza la
componente unimodular. La respuesta constitutiva completa conserva,
además, su determinante.

En particular, los lectores históricos de la forma
$\mu_r^\pm=e^{\pm2u}$ y $\varepsilon_r^\pm=e^{\mp2u}$
mantienen producto uno en cada hoja. Su razón de impedancias es
$e^{4u}$. La lectura vigente de \eqref{iii:eq:rec27-cuatro}
mantiene producto $a^4=\widehat c^{-2}$. El parámetro histórico $u$
y la coordenada $\eta$ se identifican sólo cuando se declara el mapa
entre esas realizaciones; la descomposición precedente especifica
exactamente qué factor adicional permite restituir la respuesta completa.

\subsection{Función operativa del momento de Catalán}

La restitución del teorema~\ref{iii:thm:restitucion-funcional-catalan}
sitúa la constante de Catalán en el extremo de una colección funcional.
La operación que entra en la respuesta interior es
\begin{equation}
 r_\pm=
 \frac{1+s_\pm+s_\pm^2}{s_\pm(1+s_\pm)}
 \mathcal D^2\mathcal C_2(s_\pm),\qquad
 s_\pm=q_\pm^{30},\quad \mathcal D=s\frac{d}{ds}.
 \label{iii:eq:rec27-catalan-operativo}
\end{equation}
La colección y su condición en el origen determinan el momento; la
inversión cúbica determina sus dos argumentos a partir de las respuestas.
Así se conservan dos niveles de restitución: función a partir de función
y argumentos a partir de evaluaciones en una función ya fijada.

La composición con \eqref{iii:eq:rec27-normalizacion} obtiene $a$ y
$\eta$ de esas mismas evaluaciones. La composición posterior con
\eqref{iii:eq:barbero-factorization} obtiene el funcional orientado
de Barbero--Immirzi. El valor extremo $\mathcal C_2(1)$,
las respuestas $\mathcal D^2\mathcal C_2(s_\pm)$ y el funcional
$\mathcal H(s_-)-\mathcal H(s_+)$ desempeñan, por tanto, funciones
matemáticas específicas dentro de la misma cadena. La restitución utiliza
la colección completa y las condiciones de unicidad ya demostradas.

Para precisar qué observación distingue los canales, escribamos
$u_\pm=\log r_\pm$. Las variaciones sobre una coordenatización común verifican
\begin{equation}
 \begin{pmatrix}d\log\widehat c\\d\log\widehat Z\end{pmatrix}
 =\begin{pmatrix}-1&-1\\-1&1\end{pmatrix}
 \begin{pmatrix}du_+\\du_-\end{pmatrix}.
 \label{iii:eq:rec27-variaciones}
\end{equation}
El determinante es $-2$, de modo que la lectura conjunta recupera
ambas variaciones. Cada una por separado conserva una sola combinación.
Esta identidad expresa la función complementaria de velocidad e
impedancia, con la orientación y la coordenatización dimensional fijas.

\subsection{Tres lecturas distintas de las incidencias $90/120$}

Las banderas de \eqref{iii:eq:barbero-flags} determinan las alturas
$90$ y $120$. El defecto normalizado de sus cardinales es $-1/12$.
La cadena de \eqref{iii:eq:barbero-cycle} determina, a su vez,
los pesos firmados $12$ y $-1$. Cada altura posee su serie de retorno
$S_m(q)=q^m/(1-q^{3m})$. Conservar esas tres operaciones permite
comparar de manera precisa los dos funcionales
\begin{equation}
 L_{\rm eq}(q)=S_{90}(q)-S_{120}(q),\qquad
 L_{\rm ret}(q)=12S_{90}(q)-S_{120}(q).
 \label{iii:eq:rec27-dos-kernels}
\end{equation}
El primero corresponde al antecedente de pesos iguales; el segundo es
la lectura de la cadena orientada vigente.

\begin{proposition}[Normalizaciones del polo de retorno]
Para $q\uparrow1$ se cumple
\begin{equation}
 \lim(1-q)L_{\rm eq}(q)=\frac1{1080},\qquad
 \lim(1-q)L_{\rm ret}(q)=\frac1{24},\qquad
 L_{\rm ret}-L_{\rm eq}=11S_{90}.
 \label{iii:eq:rec27-polos}
\end{equation}
\end{proposition}
\begin{proof}
La factorización de $1-q^{3m}$ prueba
$\lim(1-q)S_m(q)=1/(3m)$.
Por linealidad, los dos coeficientes son respectivamente
$1/270-1/360=1/1080$ y $12/270-1/360=1/24$.
Restar las dos definiciones prueba la última igualdad.
\end{proof}

La razón $45$ entre esos coeficientes singulares resulta de los pesos
especificados. Las funciones conservan también sus términos regulares;
la razón de los polos no identifica sus valores en todo el intervalo.
El defecto de cardinal, el peso de la cadena y el coeficiente singular
quedan así vinculados a operaciones y normalizaciones diferenciadas.

\clearpage
\subsection{Resolución de Fourier del patrón dodecafásico}

Sobre el mismo ciclo, el patrón asociado a las incidencias orientadas se
representa por
\begin{equation}
 p_m=12\cos\frac{2\pi\,3m}{12}
      -\cos\frac{2\pi\,4m}{12},\qquad 0\le m<12.
 \label{iii:eq:rec27-patron}
\end{equation}
Las frecuencias $3$ y $4$ realizan incrementos de fase $90^\circ$ y
$120^\circ$. Los argumentos de las funciones trigonométricas están
expresados en radianes. Los pesos proceden de la cadena ya construida;
la representación armónica se aplica después.

\begin{proposition}[Soporte y normalización espectral]
\label{iii:prop:rec27-fourier}
Para la transformada no normalizada
$\widehat p_k=\sum_{m=0}^{11}p_m e^{-2\pi i km/12}$,
\begin{equation}
 \widehat p_3=\widehat p_9=72,\qquad
 \widehat p_4=\widehat p_8=-6,
 \qquad \widehat p_k=0\ \text{en los demás índices}.
 \label{iii:eq:rec27-fourier}
\end{equation}
La media del patrón es cero y
$\sum_m|p_m|^2=870$.
\end{proposition}
\begin{proof}
Escribir cada coseno como la semisuma de dos exponenciales produce
coeficientes $6$ en los caracteres $3,9$ y $-1/2$ en $4,8$.
La ortogonalidad finita
$\sum_{m=0}^{11}e^{2\pi i(j-k)m/12}=12\delta_{jk}$
da \eqref{iii:eq:rec27-fourier}. El coeficiente de índice cero se
anula. La misma ortogonalidad da
$\sum|p_m|^2=(2\cdot72^2+2\cdot6^2)/12=870$.
\end{proof}

Con la transformada dividida por $12$, los coeficientes son $6$ y
$-1/2$; con la normalización unitaria son $12\sqrt3$ y $-\sqrt3$.
La razón de módulos $12:1$ se conserva en ambas. La sucesión exacta
\begin{equation}
 2(p_0,\ldots,p_{11})
 =(22,1,-23,-2,25,1,-26,1,25,-2,-23,1)
 \label{iii:eq:rec27-patron-racional}
\end{equation}
permite comprobar la transformada mediante aritmética algebraica exacta.

La media nula del patrón pertenece a la cancelación de caracteres en el
ciclo. El defecto $-1/12$ pertenece al conteo normalizado de incidencias;
$1/24$ pertenece al polo de la serie de retorno. Mantener esos dominios
explica qué conserva cada lectura antes de componerla con el par angular
y con el funcional de Barbero--Immirzi.
